%===================================================================================================
\section{Research Questions} %++++++++++++++++++++++++++++++++++++++++++++++++++++++++++++++++++++++
  \label{Research Questions} %++++++++++++++++++++++++++++++++++++++++++++++++++++++++++++++++++++++
%===================================================================================================


\improve


This thesis seeks to answer the following research questions:


%===============================================================================
\paragraph{RQ1 (Problem characterization)}
    \label{RQ1 (Problem characterization)}
Which structural characteristics of compiler and testing-framework outputs make failing
tests difficult to interpret for novice programmers in introductory programming exercises?


%===============================================================================
\paragraph{RQ2 (Design requirements)}
    \label{RQ2 (Design requirements)}
Which requirements can be derived for novice-oriented test feedback with
respect to consistency, cognitive load, and the actionability of failure information?


%===============================================================================
\paragraph{RQ3 (Reporting scheme)}
    \label{RQ3 (Reporting scheme)}
How can a uniform \emph{tested expression/expected/actual} reporting scheme be defined
such that it remains applicable across different kinds of tests (example-based unit tests
and property-based tests)?


%===============================================================================
\paragraph{RQ4 (Technical feasibility)}
    \label{RQ4 (Technical feasibility)}
How can this reporting scheme be implemented in \fs{} using Quotations and Unquote for expression
capture and evaluation, while remaining compatible with existing tooling through wrappers around
\texttt{Microsoft.\\VisualStudio.TestTools.UnitTesting} and \texttt{FsCheck}?


%===============================================================================
\paragraph{RQ5 (Results and discussion)}
    \label{RQ5 (Results and discussion)}
To what extent does the implemented artifact (\emph{Assertify} \& \emph{Checkify}) satisfy the
derived requirements in practice, and what limitations and trade-offs emerge from the wrapper-based
approach with respect to expression coverage, supported assertion forms, formatting constraints,
and integration boundaries?


By addressing these questions, this thesis positions its contribution as the design and
implementation of a test feedback layer that simplifies and standardizes failure reporting for
novice programmers, and as a critical reflection on how such a library can be extended and
integrated into teaching contexts.


%===============================================================================
% TODO +++++++++++++++++++++++++++++++++++++++++++++++++++++++++++++++++++++++++
%===============================================================================


% TODO


%===================================================================================================
% EOF ++++++++++++++++++++++++++++++++++++++++++++++++++++++++++++++++++++++++++++++++++++++++++++++
%===================================================================================================